\documentclass[10pt,a4paper]{amsart}
\usepackage[margin=19mm]{geometry}
\usepackage{amsmath,amssymb,mathtools}
\usepackage[scr=rsfs,cal=euler]{mathalfa}
\usepackage{xfrac}
\usepackage[unicode,hidelinks,bookmarks=false]{hyperref}
\newcommand{\ie}{i.e.\ }
\newcommand{\cf}{cf.\ }
\newcommand{\resp}{respectively\ }
\newcommand{\Subsection}{\subsection}
\providecommand{\nobreakdash}{\nobreak\hbox{-}}
\setlength{\emergencystretch}{2em}
\multlinegap0pt
\usepackage{bm}
\usepackage{bbm}
\usepackage{dsfont}
\usepackage{xspace}
\usepackage{cancel}
\usepackage{mathtools}
\usepackage{amsthm}
\makeatletter
\@ifundefined{theorem}{\newtheorem{theorem}{Theorem}}{}
\makeatother
\title[A formal system for reasoning about assertibility, truth, an]{A formal system for reasoning about assertibility, truth, and meaningfulness}
\author{Nik Weaver}
\date{Source: arXiv:2510.07641v1 (CC BY 4.0)}
\begin{document}
\maketitle

\noindent\emph{Source and licence.} A formal system for reasoning about assertibility, truth, and meaningfulness. Version arXiv:2510.07641v1 (CC BY 4.0), distributed under the Creative Commons Attribution 4.0 International licence, as recorded in the arXiv licence metadata for this exact version and shown on the abstract page https://arxiv.org/abs/2510.07641v1. Full rights notice: licenses/arxiv-2510.07641-CC-BY-4.0.txt.\par\smallskip

\noindent\textit{Editorial scope.} Counted unit: the Theorem whose source label is \texttt{lpcons} in the article (source0.tex lines 171--173, source label \texttt{lpcons}), delivered together with the article's own complete original proof of it (source0.tex lines 175--342, one \texttt{proof} environment of 168 lines). No mathematics is written, repaired or re-derived: statement and proof are the article's own text line for line. Required context retained uncounted: none. Every definition, display and label of those lines is retained unchanged. Omitted from this package and not redistributed: 1--170, 174--174, 343--377. Nothing inside a retained excerpt is omitted or altered: every retained body is delivered exactly as the article's lines are written, so no omission receipt of any kind accompanies this package. Citations: the retained excerpts carry no citation command at all, so no bibliography entry of the article is transcribed and no reference is redistributed. Reference adaptation: none was needed and none was made; every reference inside the delivered unit resolves inside it, so no unresolved reference or citation is produced. Printed slips kept literal: no printed word, symbol, hypothesis or number of the counted statement or of its proof was altered. Layout and class: the article's own class is \texttt{amsart} with packages including amsmath, amssymb, amsthm, enumitem; the wrapper is the lane's pinned \texttt{amsart} editorial wrapper at 10pt on A4, which restates the theorem-like environments the retained text needs and adds the article's own macro definitions that the retained text needs (none), with extra wrapper packages (bm, bbm, dsfont, xspace, cancel, mathtools). The delivered numbering is the unit's own. Assets: no retained span contains a figure environment, a graphics-inclusion command, a PDF-page inclusion, a table or a TikZ picture; no image file is copied into this package, the package declares no asset dependency and no image file is redistributed with it. Licence: version arXiv:2510.07641v1 of this article is distributed under the Creative Commons Attribution 4.0 International licence, as recorded in the arXiv licence metadata for this exact version and shown on the abstract page https://arxiv.org/abs/2510.07641v1. A full rights notice is delivered at licenses/arxiv-2510.07641-CC-BY-4.0.txt. Characters: every retained excerpt is pure ASCII, so the delivered engine, XeLaTeX with its default Latin Modern text font, reports no missing character.
\begin{theorem}\label{lpcons}
  ATM is consistent.
\end{theorem}

\begin{proof}
  We will construct a set $\mathcal{F}$ of sentences that contains $\bot$, does not contain any of the axioms, and
  whose complement is stable under the deduction rules. This shows that all the theorems of ATM lie
  in the complement of $\mathcal{F}$, and therefore $\bot$ cannot be a theorem.

  The desired set is constructed in an infinite series of {\it levels}, each of
  which consists of an infinite series of {\it stages}, each of which consists of an infinite series of {\it steps}.
  At each step new sentences are added to $\mathcal{F}$; nothing is ever removed.

  Let $S$ be any set of implications. We will use the following rules:

  \begin{itemize}
  \item ($\wedge$ rule) if $\phi \in \mathcal{F}$ and $\psi$ is any sentence, place $\phi \wedge \psi$
    and $\psi \wedge \phi$ in $\mathcal{F}$

  \item ($\vee$ rule) if $\phi, \psi \in \mathcal{F}$, place $\phi \vee \psi$ in $\mathcal{F}$

  \item ($\mathbb{T}$ rule) if $\phi \in \mathcal{F}$ then place $\mathbb{T}[t]$ in $\mathcal{F}$ for every
    term $t$ that evaluates to $\phi$.

  \item ($\to$ rule) for any implication $\phi \to \psi$ in $S$ such that  $\phi \in \mathcal{F}^c$ and
    $\psi \in \mathcal{F}$, place $\phi \to \psi$ in $\mathcal{F}$.
  \end{itemize}

  Define a function $i$ (``implication complexity'') from the set of sentences to $\omega$ by letting $i(\phi)$ be the
  number of appearances of the symbol ``$\to$'' in $\phi$, and define a function $n$ (``nesting'') from the set of
  sentences to $\omega + 1$ by setting $n(\phi) = 0$ when $\phi$ is $\bot$, $\mathbb{M}[t]$, or $\mathbb{A}[t]$, for any
  term $t$; $n(\phi \wedge \psi) = n(\phi \vee \psi) = n(\phi \to \psi) = {\rm max}(n(\phi), n(\psi))$;
  $n(\mathbb{T}[t]) = n(\hat{t}) + 1$; and $n(\phi) = \omega$ for any ungrounded sentence $\phi$.
  These functions will be used to construct $\mathcal{F}$.

  When we are at a given step of the construction I will write $\mathcal{F}_-$ for the state of $\mathcal{F}$ going
  into that step.

  The construction goes through $\omega$ levels labelled $p = 0$, $1$, $\ldots$ (``principal steps''). Every level has
  $\omega + 1$ stages labelled $n = 0$, $1$, $\ldots$, $\omega$ (corresponding to nesting). Stage $0$ on any level has $\omega$
  steps labelled $i = 0$, $1$, $\ldots$ and every other stage has $\omega$ steps labelled $i = 1$, $2$, $\ldots$ (corresponding
to implication
  complexity).

  At step 0, stage 0, level 0, when $\mathcal{F}_- = \emptyset$, we put $\bot$ in $\mathcal{F}$ and then close up
  under the $\wedge$, $\vee$, and $\mathbb{T}$ rules. That is, at the end of this step $\mathcal{F}$ is the smallest
  set that contains $\bot$ and is stable under the $\wedge$, $\vee$, and $\mathbb{T}$ rules.

  At step 0, stage 0, level 1, we place $\mathbb{M}[t]$ and $\mathbb{A}[t]$ in $\mathcal{F}$ for every term $t$
  whose evaluation $\hat{t}$ is ungrounded, and also every $\mathbb{A}[t]$ such that $\hat{t} \in \mathcal{F}_-$.
  Then we close up under the $\wedge$, $\vee$, and $\mathbb{T}$ rules.

  At step 0, stage 0 of any other level, we place every $\mathbb{A}[t]$ such that $\hat{t} \in \mathcal{F}_-$
  in $\mathcal{F}$, then close up under the $\wedge$, $\vee$, and $\mathbb{T}$ rules.

  At step $i \geq 1$, stage $n$ of any level (including $n = \omega$), let $S$ be the set of implications $\phi \to \psi$ with
  $n(\phi \to \psi) = n$ and $i(\phi \to \psi) = i$, and apply the $\to$ rule to $\mathcal{F}_-$, then
  close up under the $\wedge$, $\vee$, and $\mathbb{T}$ rules. (We do not attempt to close up under the $\to$ rule.)

  That completes the description of $\mathcal{F}$. Now we have to verify that its complement, once the construction is complete, contains the axioms of ATM and is stable under its deduction rules. Three helpful observations are

  \begin{itemize}
  \item if the sentence $\theta$ is added to $\mathcal{F}$ at step $i$ of stage $n$ on any level, then either
    $n(\theta) > n$ or else $n(\theta) = n$ and $i(\theta) \geq i$

  \item if $\phi \to \psi$ is added in to $\mathcal{F}$ anywhere on a given level, then at the end of that level
    $\phi$ will not (yet) have been added to $\mathcal{F}$

  \item at step 0, stage 0 of any level, if $\phi \not\in \mathcal{F}_-$ and $\psi \in \mathcal{F}_-$ then
    $\phi \to \psi \in \mathcal{F}_-$ (i.e., at these points in the construction $\mathcal{F}_-$ is stable under
    the $\to$ rule with $S =$ all implications).
  \end{itemize}
  The first observation is trivial when $i = n = 0$, and at any step $i \geq 1$, if $\phi \to \psi$ is added in
  to $\mathcal{F}$ then $n(\phi \to \psi) = n$ and $i(\phi \to \psi) = i$. Then we close up under the $\wedge$, $\vee$, and
  $\mathbb{T}$ rules, but the set of sentences $\theta$ with $n(\theta) > n$ or $n(\theta) = n$, $i(\theta) \geq i$ is
  stable under these rules. The second observation can be seen by noting that when $\phi \to \psi$ is added in at
  stage $n$ and step $i$, at that point $\phi$ does not belong to $\mathcal{F}$, and we also have either
  $n(\phi) < n(\phi \to \psi) = n$ or else $n(\phi) = n(\phi \to \psi) = n$ and $i(\phi) < i(\phi \to \psi) = i$.
  So by the first observation, $\phi$ has already missed any opportunity to be
  added in during the current level. The third observation holds because $n(\psi) \leq n(\phi \to \psi)$
  and $i(\psi) < i(\phi \to \psi)$, so if $\psi$ appears at any point in some level and $\phi$ does not,
  the step where $\phi \to \psi$ is
  potentially added in must come later in that level, and $\phi \to \psi$ {\it will} get added in at this later step.

  We verify that $\mathcal{F}^c$ is stable under the three deduction rules. First, the only way $\phi \wedge \psi$
  could be added is at one of the many points where we close up under the $\wedge$ rule. This means that
  if $\phi \wedge \psi$ appears in $\mathcal{F}$, then either $\phi$ or $\psi$ must also appear, showing
  that the complement of $\mathcal{F}$ is stable under the conjunction deduction rule: if neither
  $\phi$ nor $\psi$ lies in $\mathcal{F}$ then $\phi \wedge \psi$ cannot either.

  For modus ponens, suppose $\psi$ belongs to $\mathcal{F}$ and let $\phi$ be arbitrary. We want to show that
  either $\phi$ or $\phi \to \psi$ will also belong to $\mathcal{F}$. This follows from the ``second observation''
  made above: at whatever level $\psi$ first appears on, if $\phi$ has not appeared by the end of that level then
  $\phi \to \psi$ has. So if both $\phi$ and $\phi \to \psi$ are in $\mathcal{F}^c$, then so is $\psi$.

  The release rule holds because if $\psi$ appears in $\mathcal{F}$ at any level, then $\mathbb{A}[t]$ will
  appear in step 0, stage 0 of the following level, for any term $t$ that evaluates to $\psi$. Thus $\mathbb{A}[t] \in
  \mathcal{F}^c$ implies $\hat{t} \in \mathcal{F}^c$.

  Next, let us show that none of the nonlogical axioms belongs to $\mathcal{F}$. The first comment is that
  no instance of scheme (1) can appear at any point because the only place sentences of the form
  $\mathbb{M}[t]$ are placed in $\mathcal{F}$ is step 0, stage 0, level 1, and this is only done if
  $\hat{t}$ is ungrounded.
  Nor can any instance of scheme (2) be placed in $\mathcal{F}$ at any point, because if either $\hat{s}$ or
  $\hat{t}$ is ungrounded then all of $\mathbb{M}[s] \wedge \mathbb{M}[t]$, $\mathbb{M}[s \dot{\wedge} t]$,
  $\mathbb{M}[s \dot{\vee} t]$, and $\mathbb{M}[s \dot{\to} t]$ are placed $\mathcal{F}$ simultaneously, at
  step 0, stage 0, level 1; and if $\hat{s}$ and $\hat{t}$ are both grounded, none of them is ever placed in
  $\mathcal{F}$.

  Closing up under the $\wedge$, $\vee$, and $\mathbb{T}$ rules never introduces any of the axioms from schemes
  (3) through (9), because these are all implications. Nor are any of these axioms introduced at step 0 of any
  stage and level, because no implications are introduced in those steps. Thus
  we must check that no instance of any of the schemes (3) --- (9) can be added to $\mathcal{F}$ using
  the $\to$ rule at some step $i \geq 1$. Taking scheme (3), we see that the $\to$ rule will never add in
  something of the form $\mathbb{A}[t] \to \mathbb{M}[t]$, because if $\hat{t}$ is ungrounded then
  $\mathbb{A}[t]$ and $\mathbb{M}[t]$ both appear in $\mathcal{F}$ at the same time,
  and if $\hat{t}$ is grounded then $\mathbb{M}[t]$ never appears. So there is never a point
  at which $\mathbb{M}[t]$ has appeared and $\mathbb{A}[t]$ has not, which is a requirement to
  add $\mathbb{A}[t] \to \mathbb{M}[t]$ with the $\to$ rule.
  For scheme (4) observe that if either $\hat{s}$ or $\hat{t}$ is ungrounded then both sides of the implication
  appear in $\mathcal{F}$ at the same time, at step 0 of stage 0, level 1. If both are grounded then
  $\mathbb{A}[s \dot{\wedge} t]$ can only be added on step 0, stage 0 of some level such that $\hat{s}\wedge\hat{t}$
  first appeared in $\mathcal{F}$ somewhere on the previous level. But
  $\hat{s}\wedge\hat{t}$ can never be added before both $\hat{s}$ and $\hat{t}$, so at least one of
  $\mathbb{A}[s]$ and $\mathbb{A}[t]$ must be added at the same time $\mathbb{A}[s\dot{\wedge} t]$ is added,
  and thus $\mathbb{A}[s] \wedge \mathbb{A}[t]$ must also appear at the same step.

  For scheme (5), the ``second observation'' made earlier ensures that either $\hat{s}$ or $\hat{s} \to \hat{t}$
  will appear at the same level as $\hat{t}$, or earlier, so at least one of $\mathbb{A}[s]$ and
  $\mathbb{A}[s\dot{\to} t]$ (and in either case, their conjunction)
  would appear at the same step as $\mathbb{A}[t]$, at the latest.

  The left side of scheme (8) is placed in $\mathcal{F}$ at step 1, stage 0, level 0 (regardless of whether $\hat{t}$
  is grounded), whereas the right side could not be placed in $\mathcal{F}$ at any point in level 0.

  In schemes (6), (7), and (9), if $\hat{t}$ is ungrounded then both sides are placed in $\mathcal{F}$ at the same
  time, at step 0, stage 0, level 1. So assume $\hat{t}$ is grounded. Every instance of these three schemes is an
  implication with a conclusion of the form $\mathbb{A}[u]$.
  Moreover, groundedness of $\hat{t}$ entails that $\hat{u}$ is grounded in
  every case. So $\mathbb{A}[u]$ could only appear subsequent to $\hat{u}$ appearing, and thus it will suffice
  to show that $\hat{u}$ never appears. In scheme (6) we get this because groundedness of $\hat{t}$ entails that
  $\mathbb{A}[t]$ cannot appear before $\hat{t}$; in scheme (7) we never get either $\hat{t} \to \mathbb{T}[t]$
  or $\mathbb{T}[t] \to \hat{t}$, regardless of whether $\hat{t}$ is grounded, because the $\mathbb{T}$ rule is
  always invoked just after $\hat{t}$ is included, before we get a chance to apply the
  $\to$ rule to them. And scheme (9) is easy because $\mathbb{A}[t\, \dot{\leftrightarrow}\, t']$ can only be
  added after $\hat{t} \leftrightarrow \hat{t}'$ was added, and the latter can never happen, since $\hat{t} =
  \hat{t}'$.

  Finally, we must show that none of the logical axioms can appear in $\mathcal{F}$. As these are all implications,
  we only have to show that none of them can be added at any step using the $\to$ rule. This clearly cannot happen
  if any of $s$, $t$, or $u$ evaluates to an ungrounded sentence, because in that case the premise of the implication
  would have appeared at step 0, stage 0, level 1, and its conclusion could appear at that point at the earliest. If
  they all evaluate to grounded sentences, then we show that $\mathbb{A}[\dot{A}(s,t,u)]$ can never appear by checking
  that $A(\hat{s}, \hat{t}, \hat{u})$ never appears.

  This can be seen using the three observations made earlier in the proof. The verification is straightforward but tedious, so I will simply illustrate the technique for the most complicated axiom scheme,
  where $A(\hat{s}, \hat{t}, \hat{u})$ is the sentence $(\hat{s} \to (\hat{t} \to \hat{u})) \to ((\hat{s} \to \hat{t})
  \to (\hat{s} \to \hat{u}))$. If a sentence of this form
  appeared in $\mathcal{F}$ at any point,
  the second observation made above tells us that there would have to be a level at which $(\hat{s} \to \hat{t}) \to
  (\hat{s} \to \hat{u})$ is placed in $\mathcal{F}$ but $\hat{s} \to (\hat{t} \to \hat{u})$ is not. Then
  $\hat{s} \to \hat{u}$ must be placed in $\mathcal{F}$ in this level, as if it had appeared at an earlier level
  then $(\hat{s} \to \hat{t}) \to (\hat{s} \to \hat{u})$ would have too. So at the end of
  this level $\hat{u}$ is in $\mathcal{F}$ and $\hat{s}$ is not, and also $\hat{s} \to \hat{t}$ is not,
  so that since $\hat{s}$ is not in $\mathcal{F}$, $\hat{t}$ cannot be either. Thus at the end of this level we
  will have $\hat{u}$ in $\mathcal{F}$, but not $\hat{s}$ or $\hat{t}$. The second observation now tells us first that
  $\hat{t} \to \hat{u}$ must be in $\mathcal{F}$ at this point, and second that $\hat{s} \to (\hat{t} \to \hat{u})$ must
  be in $\mathcal{F}$ at this point. But this contradicts the earlier statement that
  $\hat{s} \to (\hat{t} \to \hat{u})$ does not appear in $\mathcal{F}$ at this level. We conclude that
  there is no level where $(\hat{s} \to (\hat{t} \to \hat{u})) \to ((\hat{s} \to \hat{t})
  \to (\hat{s} \to \hat{u}))$ could be added. The other logical axioms are handled similarly.
\end{proof}

\end{document}
